\subsection{The categorical reconstruction of cote 120}

\begin{saresult}{A.401}{Refuted}{Finite limits and unique witnesses do not imply images}
The reconstruction in 120/008, even with a faithful, conservative,
transportable, finite-limit-preserving functor to sets, does not imply M1.
The same failure occurs for the proposed compact-space converse in 120/018.
There is a three-element example, with every object finite, whose relation
families are closed under unique existential quantification but not under
ordinary existential quantification.
\end{saresult}
\begin{proof}
Put $R=\{a,b,c\}$, $Q=\{b,c\}$ and
$T=\{(a,b),(a,c)\}\subset R^2$. For $n\geq0$, let
$\mathcal M_n$ consist of the empty relation and all relations defined by a
finite conjunction of atoms
\[
 x_i=x_j,\qquad Q(x_i),\qquad T(x_i,x_j).
\]
The empty conjunction defines $R^n$, including $R^0=\{*\}$.
These families contain the diagonals and are closed under products,
coordinate substitutions, and finite intersections.

We verify unique projection in every arity. Merge variables identified by
equality atoms. An equivalence class with an outgoing $T$-edge must take
the value $a$; one with an incoming edge or a $Q$-atom must take a value
in $\{b,c\}$. If both occur in one class the relation is empty. Otherwise
the only restrictions, after these forced values are imposed, are the
recorded equalities. A variable class containing no retained variable and
having $Q$-type has two possible values, and a free such class has three.
They can therefore never occur in a nonempty relation whose projection
has at most one witness tuple for each retained tuple. Every eliminated
class must instead be forced to $a$ by an outgoing $T$-edge. Its target
cannot be an eliminated $Q$-class, so belongs to a retained class. Remove
the eliminated source and impose $Q$ on its retained targets. Retain all
other atoms. This is a conjunction of the same form and defines exactly
the image. Classes meeting retained variables are eliminated by equality
substitution. This proves closure under every injective coordinate
projection, with empty relations included.

Define a category $\mathcal E$ whose objects are pairs $(n,A)$ with
$A\in\mathcal M_n$. A morphism $f:(n,A)\to(m,B)$ is a function whose
graph is in $\mathcal M_{n+m}$. Identities have allowed graphs. Given
such $f:A\to B$ and $g:B\to C$, the relation
\[
 \Gamma_f(x,y)\wedge\Gamma_g(y,z)
\]
has a unique middle tuple $y$ whenever $(x,z)$ belongs to its projection;
thus the projection is $\Gamma_{gf}$ and is allowed. This proves that
composition is defined. The graph of a bijective morphism is still
allowed after interchanging its two coordinate blocks, so its inverse
is a morphism. The underlying-set functor $U$ is therefore faithful and
conservative. Products have the evident product underlying sets;
equalizers are obtained by intersecting two graphs and projecting away
the unique value. Together with $R^0$, these give finite limits preserved
by $U$. Each object embeds in a finite power of the distinguished object
$(1,R)$ by its inclusion. Replace this concrete category by its
transport-of-structure version, with objects sets equipped with a chart
to one of the displayed objects and charts identified when their
transition is an isomorphism. This equivalence makes $U$ transportable;
it neither adds relations on the fixed $R^n$ nor creates images.

For completeness, a monomorphism in $\mathcal E$ has injective underlying
function: its diagonal to its self-pullback is an isomorphism, and $U$
preserves that pullback. Conversely injective underlying maps are
monomorphisms by faithfulness. Consequently every subobject of $R^n$
has an underlying subset in $\mathcal M_n$, by the unique-projection
property just proved. Conversely each allowed subset is a subobject.
The only unary allowed sets are
\[
 \varnothing,\quad Q,\quad R.
\]
But the image of $T$ under $(x,y)\mapsto x$ is $\{a\}$.
It is not an allowed subobject, so M1 fails.

Give every underlying set the discrete topology. All objects are compact
Hausdorff, all displayed maps are continuous, and all finite limits have
the same discrete topology. The functor to compact Hausdorff spaces
has all the stated properties and the same failure of the compact
converse. The program \texttt{code/unique\_existential.py} independently
checks all relations and all coordinate projections through arity four;
the argument above, not that finite cutoff, proves closure in all arities.
\end{proof}

\begin{saresult}{A.402}{Proved under stated hypotheses}{The exact image repair}
Let $U:\mathcal E\to\mathrm{Set}$ be faithful, conservative,
transportable, and finite-limit-preserving. Let $R\in\mathcal E$ and
suppose every object embeds in some $R^n$. Identify subobjects with
their underlying subsets. Then the families
$\mathcal M_n=U(\operatorname{Sub}_{\mathcal E}(R^n))$ are closed under
unique existential quantification. They satisfy all coordinate-image
instances of M1 if and only if every morphism has a factorization
\[
 X\xrightarrow{e} I\xrightarrow{m}Y,
 \qquad U(e)\text{ surjective},\quad m\text{ monic},\quad
 U(m)(U(I))=U(f)(U(X)).
\]
Under this additional image condition, $\mathcal E$ is equivalent over
sets to the category of allowed subsets with allowed functional graphs.
It is regular and $U$ preserves regular images. Conversely regularity
with $U$ preserving images supplies exactly this condition. Thus Lead L3
is correct with this explicit repair; the hypotheses alone are insufficient.
\end{saresult}
\begin{proof}
First, monomorphisms are exactly maps sent to injections. One direction
uses the diagonal and preservation of self-pullbacks, as in A.401.
For the converse use faithfulness. If two monomorphisms into $R^n$
have the same underlying image, their pullback projections are
bijective under $U$, hence are isomorphisms by conservativity.
Thus subobjects really are determined by the specified subsets.

Products and equalizers supply intersections, coordinate substitutions,
and diagonals. If a subobject $S\hookrightarrow X\times Y$ projects
bijectively onto $X$, the projection is an isomorphism by conservativity;
it is therefore the graph of a unique morphism $X\to Y$. More generally,
the image of an injective coordinate projection of a subobject is itself
a subobject, represented by that monomorphism. This is unique-existential
closure in the sense used in A.401.

Assume M1. Embed $X$ and $Y$ into powers of $R$. The graph of
$f:X\to Y$ is a subobject of their product, so projecting it gives a
subobject $I\hookrightarrow Y$ with underlying set $U(f)(U(X))$.
The pullback $X\times_Y I\to X$ is bijective under $U$, hence an
isomorphism, and its inverse gives the required map $e:X\to I$.
Conversely the displayed factorization applied to a coordinate projection
gives every coordinate-image instance of M1.

Define the graph category using these subsets. The functor from
$\mathcal E$ sends a morphism to its graph. It is faithful, and it is
full because a functional graph has bijective first projection and this
projection is an isomorphism. It is essentially surjective by the
definition of the subsets and transportability. This proves the claimed
equivalence.

We also verify regularity, rather than assuming it silently. Let
$e:X\to I$ have surjective underlying map. If $g:X\to Z$ is constant
on the kernel pair of $e$, the image in $I\times Z$ of $(e,g)$ is the
graph of a function $h:U(I)\to U(Z)$. Its first projection is bijective,
so conservativity makes it the graph of a morphism, and $g=he$.
Surjectivity proves uniqueness. Hence $e$ is the coequalizer of its
kernel pair. Pullbacks of $e$ are still surjective under $U$ and the
same proof applies to them. This proves stable regular-epimorphism--
monomorphism factorizations. The converse follows immediately from
the assertion that $U$ takes the regular image to the set-theoretic image.
All conclusions include the empty object whenever it exists; no existence
of an empty object or of two global elements is used in this theorem.
\end{proof}
